\section*{集合论的公理与公理模式}
\phantomsection
\addcontentsline{toc}{section}{集合论的公理与公理模式}

S1. 如果 \(A\) 是一个关系，(A 或 A) \(\Longrightarrow A\) 是一条公理。

如果A和B是关系，则关系 \(A \Longrightarrow (A \text{ 或 } B)\) 是一条公理。

如果A和B是关系，则关系 \((A \text{ 或 } B) \Longrightarrow (B \text{ 或 } A)\) 是一条公理。

如果A、B和C是关系，那么关系

\[
(A \Rightarrow B) \Rightarrow ((C \text {   或   } A) \Rightarrow (C \text {   或   } B))
\]

是一条公理。

S5. 若 R 是一个关系，T 是一个项，且 x 是一个字母，则关系 \((T|x)R \Longrightarrow (\exists x)R\) 是一条公理。

S6. 设 \(\pmb{x}\) 为一个字母， \(\pmb{T}\) 并且 \(\pmb{U}\) 项，以及 \(R\{\pmb{x}\}\) 一个关系。那么该关系

\[
(T = U) \Longrightarrow (R \{T \} \Longleftrightarrow R \{U \})
\]

是一条公理。

S7. 设 R 和 S 为关系，x 为一个字母。则关系

\[
((\forall x) (R \leftrightarrow S)) \Longrightarrow (\tau_ {x} (R) = \tau_ {x} (S))
\]

是一条公理。

S8. 设 \(R\) 是一个关系， \(x\) 并且 \(y\) 不同的字母， \(X\) 并且 \(Y\) 与……不同的字母， \(x\) 并且 \(y\) 它们不出现在 \(R\)。则关系

\[
(\forall y)(\exists X)(\forall x)(R \Longrightarrow (x \in X)) \Longrightarrow (\forall Y)\operatorname{Coll}_x((\exists y)((y \in Y) \text{ 且 } R))
\]

是一条公理。

A1. \((\forall x)(\forall y)(((x\subset y)\text{ 且 }(y\subset x))\Longrightarrow (x = y)).\)

A2. \((\forall x)(\forall y)\operatorname{Coll}_z((z = x)\text{ 或 }(z = y)).\)

A3. \((\forall x)(\forall x')(\forall y)(\forall y')(((x, y) = (x', y')) \Longrightarrow ((x = x')\text{ 且 }(y = y'))).\)

A4. \((\forall X)\operatorname{Coll}_{Y}(Y \subset X).\)

A5. 存在一个无限集。
